\ifproceedings
% Proceedings title/author layout is edited in frontmatter-proceedings.tex.
% !TeX root = proceedings.tex
% Match the DOCX order and one-column title block followed by two-column text.
\vspace{.5em}
변지훈 · 유은빈 · 최서경 · 김준영 · 정소이*\par Jihun Byun · Eunbin Yu · Seokyung Choi · Junyoung Kim · Soyi Jung*\par
\vspace{.3em}
아주대학교 전자공학과\par Department of Electrical and Computer Engineering, Ajou University\par *Corresponding author: sjung@ajou.ac.kr\par

\twocolumn[{
\begin{@twocolumnfalse}
\centering
\vspace*{27bp}
{\proceedingstitlelatin\proceedingstitlehangul
 \fontsize{18bp}{28.08bp}\selectfont\bfseries
 4족 보행 로봇의 관절 토크 저하에 강건한\\
 교사, 학생 기반 강화학습 보행 제어\par}
\vspace{32bp}
{\proceedingsauthorlatin\proceedingsauthorhangul
 \fontsize{10bp}{17bp}\selectfont\bfseries
 \paperAuthorsKo\par}
{\fontsize{10bp}{17bp}\selectfont\bfseries
 \paperAffiliationKo\par}
{\fontsize{10bp}{18.12bp}\selectfont\bfseries
 \paperAuthorEmails\par}
\vspace{9bp}
{\fontsize{12bp}{18.72bp}\selectfont\bfseries
 Teacher–Student Reinforcement Learning for Robust Quadruped Locomotion under Joint Torque Degradation\par}
\vspace{19.8bp}
{\proceedingsauthorenglish
 \fontsize{10bp}{18.36bp}\selectfont\bfseries
 \paperAuthorsEn\par
 \paperAffiliationEn\par}
% Align the Abstract title and first right-column text at their prior midpoint.
\vspace{35.604bp}
\end{@twocolumnfalse}
}]
\vspace*{-16.204bp}
{\centering\fontsize{12bp}{18.72bp}\selectfont\bfseries Abstract\par}
\vspace{14.04bp}
\paperbodyfont
\indent% !TeX root = proceedings.tex
We propose a reinforcement learning locomotion policy that maintains balance, tracks commanded motion, and follows reference paths under nominal and single-joint torque degradation conditions. The policy combines current observations and history-based features through teacher–student joint training, without direct fault inputs. Policies trained with one seed each are compared in Isaac Gym using a Unitree A1 quadruped model on matched mixed-terrain tasks. Under full torque degradation, survival improved by 37.2 percentage points and mean pre-termination cross-track root mean square error decreased by 64.6\% relative to the nominally trained policy.
\par
% The proceedings template has no separate Korean summary or Keywords block.
% The original Korean summary and Keywords remain in frontmatter.tex for review.

\else
\twocolumn[{
\begin{@twocolumnfalse}
\begin{center}
{\ifproceedings\fontsize{18pt}{21.6pt}\else\fontsize{20pt}{32pt}\fi\selectfont\bfseries 4족 보행 로봇의 액추에이터 토크 저하에 강건한 교사·학생 기반 강화학습 보행 제어\par}
\vspace{.4em}
{\ifproceedings\fontsize{12pt}{14.4pt}\else\fontsize{17pt}{22.1pt}\fi\selectfont Teacher–Student Reinforcement Learning for Robust Quadruped Locomotion under Actuator Torque Degradation\par}
\ifproceedings\vspace{.5em}
변지훈 · 유은빈 · 최서경 · 김준영 · 정소이*\par Jihun Byun · Eunbin Yu · Seokyung Choi · Junyoung Kim · Soyi Jung*\par
\vspace{.3em}
아주대학교 전자공학과\par Department of Electrical and Computer Engineering, Ajou University\par *Corresponding author: sjung@ajou.ac.kr\par
\fi
\end{center}
{\fontsize{9pt}{11.7pt}\selectfont
\ifproceedings\else\leftskip=17pt\rightskip=17pt\fi
\noindent\textbf{요약}\quad 단일 관절의 토크 저하에 대응하는 기존 교사·학생 공동학습 방식을 적용한 4족 보행 정책을 구성한다. 특권 물리 정보를 사용하는 교사와 관측 이력을 사용하는 학생을 공동학습하고, 평가 시에는 고장 정보를 직접 받지 않는 학생만 사용한다. Isaac Gym의 Unitree A1에서 12개 관절과 6개 손실 수준을 시험하였다. 완전 손실에서 생존율은 정상 학습 정책의 41.49~\%에서 88.89~\%로 증가하였다. 전진속도 RMSE는 0.317 m/s에서 0.175 m/s로, yaw-rate RMSE는 0.753 rad/s에서 0.268 rad/s로 감소하였다. 관절별 분석으로 고장 후 명령 추종의 취약 조건을 구분하였다.\par\ifproceedings\medskip\else\vspace{16pt}\fi
\noindent\textbf{Abstract}\quad We construct a quadruped policy for single-joint torque loss by applying an existing joint-training framework with privileged Teacher and history-based Student representations. Inference uses only the Student without explicit fault information. We evaluate Unitree A1 in Isaac Gym across 12 joints and six loss settings. Under complete torque loss, survival increases from 41.49~\% for a fault-naive policy to 88.89~\%. Forward-velocity RMSE decreases from 0.317 to 0.175~m/s, and yaw-rate RMSE decreases from 0.753 to 0.268~rad/s. Joint-wise analysis identifies conditions with reduced command-tracking performance after fault onset.\par
\smallskip\noindent\textbf{Keywords:} Quadrupedal locomotion, Actuator degradation, Teacher–student learning, Joint learning
}
\vspace{1em}
\end{@twocolumnfalse}
}]

\fi
